\documentclass[11pt,a4paper]{article}
\usepackage[margin=27mm]{geometry}
\usepackage[T1]{fontenc}
\usepackage{lmodern,amsmath,amssymb,amsthm,microtype}
\usepackage[hidelinks]{hyperref}
\hypersetup{pdftitle={Computer verification for "Explicit graphs with quantum chromatic number three and large chromatic number"},
  pdfauthor={Julius A. Zeiss},pdfdisplaydoctitle=true}

\newtheorem{lemma}{Lemma}

\setlength{\emergencystretch}{2em}
\newcommand{\F}{\mathbb{F}_2}
\newcommand{\RR}{\mathbb{R}}

\title{Computer verification for\\ ``Explicit graphs with quantum chromatic
number three\\ and large chromatic number''}
\author{Julius A. Zeiss\\[4pt]
{\small Institute for Quantum Information, RWTH Aachen University}\\
{\small Otto-Blumenthal-Stra{\ss}e 20, 52074 Aachen, Germany}\\
{\small \href{mailto:jzeiss@physik.rwth-aachen.de}{\texttt{jzeiss@physik.rwth-aachen.de}}}}
\date{19 November 2026}

\begin{document}
\maketitle

This note accompanies the paper~\cite{Z}. The proofs of the theorems,
propositions and lemmas of~\cite{Z} do not rely on computation. The programs
in this repository check its finite statements independently, and this note
describes what they verify. The numbers of sections, statements and
equations below are those of~\cite{Z}, and so are the notation and the
terms. In particular $F$ is the graph on $219$ vertices
of~\cite[Section~3]{Z}, the graphs on $28$, $29$, $32$ and $165$ vertices are
those of~\cite[Section~6 and Appendix~C]{Z}, and $\chi$ is the chromatic
number.

Both programs need Python 3.8 or later and use only the standard library.
All arithmetic is exact integer arithmetic. No numerical solver is involved.

\section{A chain certificate for the graph on 219 vertices}

For a finite graph $G$, let $A(G)$ be the box complex
$\mathsf{B}_{\mathrm{edge}}(G)$ of Matou\v{s}ek and Ziegler~\cite{MZ}. Its
vertices are the directed edges $(a,b)$ of $G$, which we call \emph{arcs}.
A set of arcs is a simplex
when $a\sim b'$ for every two of its members $(a,b),(a',b')$, including
an arc paired with itself. Reversal $T(a,b)=(b,a)$ acts freely on
$|A(G)|$, since the first and second endpoints of a simplex form
disjoint sets. A proper $m$-coloring $f$ gives an affine map on each
simplex by $(a,b)\mapsto e_{f(a)}-e_{f(b)}$. The colors on its first
and second endpoints are disjoint, so this map is nowhere zero.
It takes values in the coordinate-sum-zero hyperplane of $\RR^m$ and
changes sign under $T$. Normalizing gives a map to $S^{m-2}$ that commutes
with $T$ and the antipodal map. As in~\cite{Z} we call such a map
\emph{odd}.

\begin{lemma}\label{lem:certificate}
Suppose $c_0,\ldots,c_k$ are simplicial chains of $A(G)$ over $\F$,
with $c_i$ of dimension $i$, $c_0$ a single vertex, and
$\partial c_i=c_{i-1}+Tc_{i-1}$ for $1\le i\le k$, where $k\ge1$.
Then $|A(G)|$ admits no odd map to $S^{k-1}$ and $\chi(G)\ge k+2$.
\end{lemma}

\begin{proof}
Let $p:|A(G)|\to Q=|A(G)|/T$ be the orbit map and $w\in H^1(Q;\F)$
the first Stiefel--Whitney class of this double cover. Each quotient cell
has two lifts. Over $\F$,
summing these lifts and projecting to the quotient give an exact
sequence of cellular chain complexes
\[
 0\longrightarrow C_*(Q)\longrightarrow C_*(A(G))
 \xrightarrow{p_\#} C_*(Q)\longrightarrow0 .
\]
Its connecting homomorphism is cap product with $w$. This is the homology
Gysin sequence of the double cover $p$, the analogue for a $0$-sphere
bundle of the sequence~(6) of~\cite{Z}; see
\cite[Chapter~5, \S 7]{Sp}, \cite[\S 12]{MS}. Each $p_\#c_i$ is a cycle and
the boundary identities give $w\cap[p_\#c_i]=[p_\#c_{i-1}]$.
Hence $w^k\cap[p_\#c_k]$ is the nonzero class of a point.
An odd map to $S^{k-1}$ would pull $w$ back from the target
quotient $\RR P^{k-1}$, where every degree-$k$ class vanishes.
The coloring map above now gives $\chi(G)\ge k+2$.
\end{proof}

The certificate lists chains $c_0,c_1,c_2,c_3$ in $A(F)$
of respective sizes $(1,3,684,5\,832)$. The verifier reconstructs
the induced graph $F$ from the labelings of~\cite[Section~3]{Z} and
checks its $1\,353$ edges. For each chain it checks every cross-adjacency
in every simplex and forms its boundary by cancellation modulo two. The
three boundary equalities give $w^3\ne0$ for the double cover
$|A(F)|\to |A(F)|/T$, by the proof of Lemma~\ref{lem:certificate}, and
therefore exclude a four-coloring. This is a second proof of
$\chi(F)\ge5$, independent of~\cite[Corollary~3.5]{Z}.

The data file \texttt{graph\_219\_certificate.json.gz} lists
the vertices, edges, proper five-coloring, and chains.
From the directory of the repository run
\begin{center}
  \texttt{python3 verify\_graph\_219.py}
\end{center}
The verifier also checks the identities for the projections, commutation
on every edge, and the three constant vertices.

The same boundary identities have been verified for the glued graphs
$G_1(C_n)$ with $n=3,5,7,9$, and for the wired graphs $G_1^\lambda(C_n)$
with
\[
 (n,\lambda)\in\{(3,0),(3,1),(3,3),(5,0),(5,2),(7,1),(9,0)\}.
\]
For $G_1(C_5)$ the top chain has $27\,000$ simplices supported on
$1\,085$ graph vertices. These eleven computations give independent
checks of the analytical first-level lower bound
of~\cite[Theorem~1.1]{Z}. The chain data in the repository are those for
$F$.

\section{Verification of the small graphs}

The classical and quantum chromatic numbers of the graphs on $28$,
$29$, $32$ and $165$ vertices and of the graph $F$ on $219$ vertices have
analytical proofs in~\cite{Z}. Theorem~6.4 and Proposition~C.3 of~\cite{Z}
prove vertex-criticality on $28$, $165$ and $29$ vertices by hand. The
$32$-vertex parent and its $29$-vertex induced subgraph have the same
chromatic number \cite[Remark~C.4]{Z}.

The mod~$2$ chain certificate of the previous section proves the stronger
topological statement $w^3\ne0$ for the double cover
$|A(F)|\to |A(F)|/T$ of the graph $F$ \cite[Proposition~3.6]{Z}. It
gives a second proof of $\chi(F)\ge5$ alongside the elementary proof
in~\cite[Corollary~3.5]{Z}. The eleven additional chain computations
reported above check analytical lower bounds for other graphs. The lower
bounds on the minimum order of a separation in~\cite[Section~7]{Z} use
Lalonde's computer-assisted theorem~\cite{La}.

For the small graphs run
\begin{center}
  \texttt{python3 verify\_small\_graphs.py}
\end{center}
Its data file is \texttt{small\_graphs\_data.json.gz}. The program
reconstructs the graphs on $28$, $29$, $32$ and
$165$ vertices from their measurements, counts their vertices and edges,
and checks the relations~(8) of~\cite{Z}. As independent checks of
Theorem~6.4 and Proposition~C.3, it searches for three-colorings
after vertex deletions on $28$ and $29$ vertices and tests the $165$
stored four-colorings on the corresponding vertex-deleted subgraphs.
It also checks proper colorings by witness vectors, in the manner of the
proof of~\cite[Corollary~C.5]{Z}, with four colors on $28$ and
$32$ vertices and five colors on $165$ vertices, and the local
color tables for the $28$- and $29$-vertex graphs
in~\cite[Proposition~C.2]{Z}, using the original order of each basis.

The program independently checks the parity rule
\cite[Proposition~6.2]{Z}. Each of the three plane grids with its
constants has $44$ normalized colorings, $22$ of each truth value. Cells of
different planes are nonadjacent. Of the $44^3=85\,184$ colorings of the
three grids, $42\,592$ have even parity. The colored neighbors already
exclude all four colors at some remaining vertex in $18\,804$ cases. In
each of the other $23\,788$ cases, repeatedly assigning a vertex its sole
available color produces an empty list.

As checks of Lemmas~3.1--3.3 of~\cite{Z}, the program also enumerates
$720$ normalized colorings of the two grids of orthogonal,
noncomplementary projections used in the check. Six make both projections
true. Complementary projections have the same grid and opposite truth
values. The gadget of a three-part measurement has $7\,986$ normalized
colorings. Finally the program compares the measurements
in~\cite[Corollary~3.5]{Z} with the vertex list of the certificate of
$F$, checks its $1\,353$ edges, and counts the $111$ other orthogonal
pairs on the same vertices. For the $597$ measurements in
$\mathcal M^{(1)}$ \cite[Section~4.5]{Z}, it counts $6\,024$ orthogonal
pairs, of which $5\,538$ occur within a common context.

\begin{thebibliography}{9}

\bibitem{La}
O.~Lalonde,
\emph{On the quantum chromatic numbers of small graphs},
The Electronic Journal of Combinatorics \textbf{32} (2025), no.~1, P1.18.
\href{https://arxiv.org/abs/2311.08194}{arXiv:2311.08194}.

\bibitem{MZ}
J.~Matou\v{s}ek and G.~M.~Ziegler,
\emph{Topological lower bounds for the chromatic number: A hierarchy},
Jahresbericht der Deutschen Mathematiker-Vereinigung \textbf{106} (2004),
71--90.
\href{https://arxiv.org/abs/math/0208072}{arXiv:math/0208072}.

\bibitem{MS}
J.~W.~Milnor and J.~D.~Stasheff,
\emph{Characteristic Classes},
Annals of Mathematics Studies 76, Princeton University Press, 1974.

\bibitem{Sp}
E.~H.~Spanier,
\emph{Algebraic Topology},
McGraw-Hill, 1966.

\bibitem{Z}
J.~A.~Zeiss,
\emph{Explicit graphs with quantum chromatic number three and large
chromatic number}, 2026.

\end{thebibliography}
\end{document}
